\chapter{面積・体積・長さと運動}\label{章:積分の応用}

定積分は，細かく分けた量を足し合わせる極限として定義した。この考えを使うと，曲線に囲まれた面積だけでなく，立体の体積や曲線の長さも求められる。ただし，どの問題にも被積分関数が最初から与えられるわけではない。細い部分を一つ取り出し，その面積，体積，長さをどの式で近似するかを決める必要がある。

その際に確かめたいことが三つある。一つは，分割した部分が対象全体を覆っているか。もう一つは，同じ部分を重複して数えていないか。そして，面積や道のりのような非負の量に，符号付きの差をそのまま使っていないかである。この章では，具体的な図形と運動を通して，それぞれの確認がどの式に現れるかを見ていく。

\section{曲線で囲まれた面積}\label{節:曲線間の面積}

$a<b$ とし，$f,g$ は $[a,b]$ で連続で，$f(x)\geq g(x)$ を満たすとする。二曲線 $y=f(x)$，$y=g(x)$ の間を，$x$ 軸に垂直な細い帯に分ける。幅が $\Delta x$ の帯の高さは，その付近での $f(x)-g(x)$ である。したがって面積は
\begin{equation}\label{式:上下の差の面積}
S=\int_a^b\{f(x)-g(x)\}\dd x
\end{equation}
になる。上下の曲線がともに $x$ 軸の下にあっても，高さは上の値から下の値を引いた差であり，同じ式が使える。

二曲線が途中で交差するなら，上下関係が変わる点で区間を分ける。一つの式で表せば
\[
S=\int_a^b|f(x)-g(x)|\dd x
\]
であるが，実際に計算するには差の符号を調べて絶対値を外す。交点の $x$ 座標は区切りの候補であり，交点を求めただけで上下が決まるわけではない。差を因数分解したり，各区間で一つの値を代入したりして確かめる。

\begin{example}{途中で上下が入れ替わる面積}
二曲線 $y=x^3-3x$ と $y=x$ に囲まれた部分の面積を求めよ。
\end{example}

交点は $x^3-3x=x$，すなわち $x(x-2)(x+2)=0$ より，$x=-2,0,2$ である。差 $x^3-4x$ は $-2<x<0$ で正，$0<x<2$ で負なので，面積は
\[
S=\int_{-2}^0(x^3-4x)\dd x
 +\int_0^2(4x-x^3)\dd x.
\]
図\ref{図:面積の上下交換}では，左右の部分が原点について対称である。積分でも第一項に $x=-t$ を使えば第二項と同じになるから，
\[
S=2\int_0^2(4x-x^3)\dd x
 =2\left[2x^2-\frac{x^4}{4}\right]_0^2=8.
\]
これを $\int_{-2}^2(x^3-4x)\dd x$ とすれば奇関数の積分なので零になるが，その値は左右の面積を一方は正，一方は負として足したものである。

\begin{figure}[htbp]
\centering
\begin{tikzpicture}[x=1.6cm,y=0.65cm,>=Stealth]
\fill[blue!12]plot[domain=-2:0,samples=60](\x,{\x*\x*\x-3*\x})--plot[domain=0:-2,samples=30](\x,\x)--cycle;
\fill[orange!17]plot[domain=0:2,samples=30](\x,\x)--plot[domain=2:0,samples=60](\x,{\x*\x*\x-3*\x})--cycle;
\draw[->](-2.3,0)--(2.4,0)node[right]{$x$};
\draw[->](0,-2.8)--(0,3.4)node[above]{$y$};
\draw[blue!70!black,thick]plot[domain=-2.13:2.13,samples=120](\x,{\x*\x*\x-3*\x});
\draw[orange!80!black,thick](-2.2,-2.2)--(2.2,2.2);
\foreach \x in {-2,2}{\draw[dashed](\x,0)--(\x,\x);\node[below]at(\x,0){$\x$};}
\node[below right]at(0,0){$O$};
\node[blue!70!black]at(-1.25,2.8){$y=x^3-3x$};
\node[orange!80!black]at(1.6,3.0){$y=x$};
\end{tikzpicture}
\caption{同じ二曲線でも，交点の左右で上側の曲線が入れ替わる。}\label{図:面積の上下交換}
\end{figure}

\subsection{横に切るほうが簡単な場合}

細い帯は縦に取る必要はない。$c\leq y\leq d$ において，右側の境界が $x=q(y)$，左側が $x=p(y)$ と表せるなら，帯の横の長さは $q(y)-p(y)$ なので
\begin{equation}\label{式:左右の差の面積}
S=\int_c^d\{q(y)-p(y)\}\dd y
\end{equation}
である。式\eqref{式:上下の差の面積}では上下の差，この式では左右の差を使う。変数の文字を交換するだけでなく，実際にどの向きに帯を取るかが変わっている。

たとえば $x=y^2$ と $x=2-y$ に囲まれる部分では，交点の $y$ 座標は $-2,1$ である。その間では直線が右側なので，
\[
S=\int_{-2}^1(2-y-y^2)\dd y
 =\left[2y-\frac{y^2}{2}-\frac{y^3}{3}\right]_{-2}^1
 =\frac92.
\]
縦に切ると放物線の上側と下側を $y=\pm\sqrt x$ に分け，さらに直線との関係も場合分けすることになる。横に切れば，一つの区間で一つの差になる。境界をどちらの変数で表すと少ない場合分けで済むかが，積分変数を選ぶ基準になる。

\section{断面積を積み重ねて体積を求める}\label{節:断面積と体積}

立体を $x$ 軸に垂直な平面で切り，位置 $x$ における断面積を $A(x)$ とする。幅 $\Delta x$ の薄い部分を，底面積 $A(x)$，高さ $\Delta x$ の柱体で近似すれば，その体積は $A(x)\Delta x$ である。断面積が連続に変化する通常の立体では，分割を細かくした極限から
\begin{equation}\label{式:断面積の体積}
V=\int_a^b A(x)\dd x
\end{equation}
を得る。断面が途中で形を変える場合も，変わる位置で有限個の区間に分け，各区間で同じ考えを使える。

この式の $A(x)$ は長さではなく面積である。長さの単位を $\mathrm{cm}$ とすれば，$A(x)$ は $\mathrm{cm}^2$，幅 $\dd x$ を掛けた体積は $\mathrm{cm}^3$ になる。図形の式が合っているか迷ったとき，単位も確認の助けになる。

\begin{example}{正方形の断面を持つ立体}
平面上で $0\leq x\leq1$，$x^2\leq y\leq x$ となる部分を底面とする。各 $x$ において，底面内の線分を一辺とする正方形が，底面に垂直な断面になるような立体の体積を求めよ。
\end{example}

底面内の線分の長さは $x-x^2$ である。この長さをそのまま積分すると底面積になるが，今求めたいのは体積なので，正方形の断面積
\[
A(x)=(x-x^2)^2=x^2-2x^3+x^4
\]
を積分する。
\[
V=\int_0^1(x^2-2x^3+x^4)\dd x
 =\left[\frac{x^3}{3}-\frac{x^4}{2}+\frac{x^5}{5}\right]_0^1
 =\frac1{30}.
\]
同じ底面でも，断面が正三角形なら面積は $\frac{\sqrt3}{4}(x-x^2)^2$ であり，体積も上の値の $\frac{\sqrt3}{4}$ 倍になる。底面の情報だけでは体積は決まらず，断面の形が必要である。

\section{回転体では軸からの距離を測る}\label{節:回転体の体積}

$0\leq g(x)\leq f(x)$ の領域 $g(x)\leq y\leq f(x)$ を $x$ 軸のまわりに回すと，位置 $x$ の断面は外半径 $f(x)$，内半径 $g(x)$ の円環になる。よって
\begin{equation}\label{式:円環断面の体積}
V=\pi\int_a^b\{f(x)^2-g(x)^2\}\dd x.
\end{equation}
下側が軸そのもの，すなわち $g(x)=0$ なら断面は円板で，$V=\pi\int_a^b f(x)^2\dd x$ となる。

回転軸が $y=c$ なら，半径は $y$ 座標そのものではなく $|y-c|$ である。たとえば $c\leq g(x)\leq f(x)$ なら，外半径は $f(x)-c$，内半径は $g(x)-c$ になる。軸が領域より上にあるときは，下側の境界のほうが軸から遠い。上下の順序をそのまま外半径・内半径の順序にしてはいけない。

\subsection{軸を横切る断面と，回転後の重なり}

縦の線分 $g(x)\leq y\leq f(x)$ が回転軸 $y=c$ を横切ると，回転後の断面には中心まで点が入る。したがって内半径は零であり，外半径は両端の遠いほうである。まとめると，
\[
R(x)=\max\{|f(x)-c|,|g(x)-c|\},
\]
\[
r(x)=
\begin{cases}
0,&g(x)\leq c\leq f(x),\\
\min\{|f(x)-c|,|g(x)-c|\},&\text{それ以外}
\end{cases}
\]
として，断面積は $\pi\{R(x)^2-r(x)^2\}$ になる。

軸より上と下を別々に回すと，同じ円板の部分を二度作ることがある。その場合，二つの体積を加えると重複して数える。先に回転後の断面を考え，できた円板や円環の和集合の面積を求めることが必要である。

\begin{example}{軸をまたぐ領域を回す}
$0\leq x\leq2$，$x-1\leq y\leq1$ で表される領域を $x$ 軸のまわりに一回転してできる立体の体積を求めよ。
\end{example}

$0\leq x\leq1$ では線分が軸をまたぐ。上端の距離は $1$，下端の距離は $1-x\leq1$ なので，断面は半径 $1$ の円板になる。$1\leq x\leq2$ では線分全体が軸の上にあり，外半径 $1$，内半径 $x-1$ の円環になる。図\ref{図:回転断面の区別}は，この違いを断面で示している。
\[
\begin{aligned}
V&=\pi\int_0^1 1\dd x
  +\pi\int_1^2\{1-(x-1)^2\}\dd x\\
 &=\pi+\pi\left[x-\frac{(x-1)^3}{3}\right]_1^2
 =\frac{5\pi}{3}.
\end{aligned}
\]
最初から $\pi\int_0^2\{1-(x-1)^2\}\dd x$ とすると，$x<1$ でも中心の円板をくり抜いてしまう。しかし元の線分は軸を含んでおり，回転後にその穴は存在しない。

\begin{figure}[htbp]
\centering
\begin{tikzpicture}[scale=1.0,>=Stealth]
\begin{scope}[xshift=-3.6cm]
\fill[blue!12](0,-1)--(0,1)--(2,1)--cycle;
\draw[->](-0.3,0)--(2.5,0)node[right]{$x$};
\draw[->](0,-1.3)--(0,1.5)node[above]{$y$};
\draw[thick](0,1)--(2,1)--(0,-1);
\draw[blue!70!black,very thick](0.5,-0.5)--(0.5,1);
\draw[orange!85!black,very thick](1.5,0.5)--(1.5,1);
\node[below]at(1,0){$1$};\node[below]at(2,0){$2$};
\node[below]at(1,-1.7){回す前の領域};
\end{scope}
\begin{scope}[xshift=0.1cm]
\fill[blue!12](0,0)circle(1);
\draw[blue!70!black,thick](0,0)circle(1);
\draw[->](0,0)--(0.707,0.707)node[midway,left]{$1$};
\fill(0,0)circle(0.025);
\node[below]at(0,-1.45){$x=\frac12$：円板};
\end{scope}
\begin{scope}[xshift=3.2cm]
\fill[orange!17,even odd rule](0,0)circle(1)(0,0)circle(0.5);
\draw[orange!85!black,thick](0,0)circle(1);
\draw[orange!85!black,thick](0,0)circle(0.5);
\draw[->](0,0)--(0.707,0.707)node[midway,right]{$1$};
\draw[->](0,0)--(-0.5,0)node[midway,above]{$\frac12$};
\fill(0,0)circle(0.025);
\node[below]at(0,-1.45){$x=\frac32$：円環};
\end{scope}
\end{tikzpicture}
\caption{軸を含む線分は円板を作り，軸から離れた線分は円環を作る。}\label{図:回転断面の区別}
\end{figure}

一つの $x$ に対して元の領域が複数の線分に分かれる場合には，各線分が作る半径の範囲を調べる。半径の範囲が重なるなら，その部分を一度だけ数えるように合わせる。この確認は，元の平面図形では重なっていない二つの部分を回す場合にも必要になる。

\subsection{円筒の薄い層として分ける}

同じ立体でも，回転軸に垂直に切る代わりに，軸に平行な薄い帯を回す方法がある。たとえば $0\leq a\leq x\leq b$，$g(x)\leq y\leq f(x)$ を $y$ 軸のまわりに回す。幅 $\Delta x$ の帯は，半径がほぼ $x$，高さがほぼ $f(x)-g(x)$ の薄い円筒の層になる。

半径 $x$ と $x+\Delta x$ の間の円環の面積は
\[
\pi\{(x+\Delta x)^2-x^2\}
 =2\pi x\Delta x+\pi(\Delta x)^2
\]
である。分割を細かくすると二次の項による誤差の和は零に近づくので，体積は
\begin{equation}\label{式:円筒殻の体積}
V=2\pi\int_a^b x\{f(x)-g(x)\}\dd x
\end{equation}
となる。これを円筒殻による方法という。$2\pi x$ は円周の長さ，$f-g$ は層の高さである。

\begin{example}{同じ体積を二通りに表す}
$0\leq x\leq1$，$x^2\leq y\leq x$ の領域を $y$ 軸のまわりに回転させた立体の体積を求めよ。
\end{example}

縦の帯を回すと，半径 $x$，高さ $x-x^2$ の円筒殻になるから，
\[
V=2\pi\int_0^1x(x-x^2)\dd x
 =2\pi\left[\frac{x^3}{3}-\frac{x^4}{4}\right]_0^1
 =\frac{\pi}{6}.
\]
横に切る方法でも確かめられる。$0\leq y\leq1$ で，領域の左右の端は $x=y$ と $x=\sqrt y$ である。したがって円環の外半径は $\sqrt y$，内半径は $y$ となり，
\[
V=\pi\int_0^1\{(\sqrt y)^2-y^2\}\dd y
 =\pi\left[\frac{y^2}{2}-\frac{y^3}{3}\right]_0^1
 =\frac{\pi}{6}.
\]
断面を円環として積み重ねるか，円筒の層として重ねるかによって，被積分関数は異なる。しかし，それぞれが立体を重複なく分けているので同じ体積になる。

式\eqref{式:円筒殻の体積}で $x\geq0$ としたのは，異なる $x$ が異なる半径に対応するためである。領域が $y$ 軸をまたぐと，$x$ と $-x$ が同じ半径になる。たとえば $-1\leq x\leq1$，$0\leq y\leq1$ の長方形を回すと半径 $1$，高さ $1$ の円柱で，体積は $\pi$ である。ところが $2\pi\int_{-1}^1|x|\dd x=2\pi$ は同じ層を二度数えている。絶対値で半径を正にしただけでは，重複は解消しない。

\section{曲線の長さと媒介変数}\label{節:曲線の弧長}

曲線 $y=f(x)$ の長さを求めるとき，細い部分を長方形で近似するのは適切ではない。曲線上の近い二点を線分で結び，折れ線の長さで近似する。$f$ が $[a,b]$ で連続微分可能，すなわち導関数も連続であるとする。小区間の幅を $\Delta x$，両端の関数値の差を $\Delta y$ とすると，線分の長さは
\[
\sqrt{(\Delta x)^2+(\Delta y)^2}
 =\sqrt{1+\left(\frac{\Delta y}{\Delta x}\right)^2}\,\Delta x.
\]
平均値の定理により，小区間内のある $\xi$ で $\frac{\Delta y}{\Delta x}=f'(\xi)$ である。したがって折れ線の長さの和は，連続関数 $\sqrt{1+f'(x)^2}$ のリーマン和になり，分割を細かくすると
\begin{equation}\label{式:グラフの弧長}
L=\int_a^b\sqrt{1+f'(x)^2}\dd x
\end{equation}
に近づく。これが曲線の弧長である。有限個の点で折れ曲がっていても，各部分が連続微分可能なら，部分ごとの長さを加えればよい。

$\int_a^b|f'(x)|\dd x$ は縦方向の変化の総量であり，弧長ではない。水平な線分では $f'=0$ だが長さは $b-a$ である。式\eqref{式:グラフの弧長}の平方根の中の $1$ は，横方向の移動も数えている。

\subsection{媒介変数で表された曲線}

円や，縦向きの接線を持つ曲線は，一つの関数 $y=f(x)$ で扱いにくい場合がある。そこで
\[
x=x(t),\qquad y=y(t),\qquad \alpha\leq t\leq\beta
\]
と表す。$t$ を少し増やしたとき，横の変化はほぼ $x'(t)\Delta t$，縦の変化はほぼ $y'(t)\Delta t$ だから，線分の長さはほぼ
\[
\sqrt{x'(t)^2+y'(t)^2}\,\Delta t
\]
である。$x,y$ が連続微分可能なら，導関数の連続性によりこの近似の誤差は区間全体で制御でき，折れ線の極限として
\begin{equation}\label{式:媒介曲線の弧長}
L=\int_\alpha^\beta\sqrt{x'(t)^2+y'(t)^2}\dd t
\end{equation}
を得る。ここでは近似誤差の評価の詳細は省き，この弧長公式を用いる。途中で速度に当たる平方根が零になる点があっても，この式自体は使える。

ただし，この積分は $t$ に沿って進んだ長さを数える。同じ線分や曲線を二度たどれば二度数える。たとえば $x=\cos t$，$y=\sin t$ を $0\leq t\leq4\pi$ で積分すると $4\pi$ だが，点の集合としての単位円の周長は $2\pi$ である。図形そのものの長さを求めるなら，一度だけたどる区間を選ぶ。

媒介変数は面積にも使える。$x(t)$ が単調増加し，$y(t)\geq0$ で，曲線が $x$ の関数として表される場合は，置換 $\dd x=x'(t)\dd t$ により
\begin{equation}\label{式:媒介曲線の面積}
S=\int_\alpha^\beta y(t)x'(t)\dd t
\end{equation}
となる。$x(t)$ が途中で折り返すなら，区間を分けて曲線が囲む部分を確認する。この式に絶対値を付けるだけで常に図形の面積になるわけではない。帯が重なれば，同じ面積を二度数えるからである。

\begin{example}{一つの媒介表示から面積と長さを求める}
$a>0$ とする。曲線
\[
x=a(t-\sin t),\qquad y=a(1-\cos t),\qquad 0\leq t\leq2\pi
\]
の長さと，この曲線と $x$ 軸に囲まれた部分の面積を求めよ。
\end{example}

$x'(t)=a(1-\cos t)$，$y'(t)=a\sin t$ である。まず弧長の平方根を計算すると
\[
\begin{aligned}
x'(t)^2+y'(t)^2
 &=a^2\{(1-\cos t)^2+\sin^2t\}\\
 &=2a^2(1-\cos t)=4a^2\sin^2\frac t2.
\end{aligned}
\]
平方根は $2a|\sin\frac t2|$ である。この区間では $0\leq\frac t2\leq\pi$ なので正弦は非負であり，
\[
L=2a\int_0^{2\pi}\sin\frac t2\dd t
 =\left[-4a\cos\frac t2\right]_0^{2\pi}=8a.
\]
平方根を取る際に絶対値が必要かどうかを，変数の範囲で確認している。

一方，$x'$ は端点以外で正なので，$x$ は区間全体で単調増加する。また $y\geq0$ で，両端では $y=0$ である。よって面積は
\[
\begin{aligned}
S&=a^2\int_0^{2\pi}(1-\cos t)^2\dd t\\
 &=a^2\int_0^{2\pi}\left(\frac32-2\cos t+\frac12\cos2t\right)\dd t
 =3\pi a^2.
\end{aligned}
\]
弧長では二方向の変化から距離を作り，面積では高さに横の幅を掛けた。同じ媒介表示でも，求める量に応じて積分の中身が異なる。

\section{速度から位置と道のりへ}\label{節:運動と積分}

一直線上の運動では，位置を $s(t)$，速度を $v(t)=s'(t)$，加速度を $a(t)=v'(t)$ とする。速度が連続なら，微積分の基本定理により
\[
s(t_1)-s(t_0)=\int_{t_0}^{t_1}v(t)\dd t.
\]
左辺は位置の差，すなわち変位である。出発点に戻れば零になるが，途中で移動した距離が零という意味ではない。

速度が正なら正の向きに進み，負なら負の向きに進む。移動した距離を向きによらず加えるには速さ $|v(t)|$ を積分して，
\begin{equation}\label{式:道のり}
\text{道のり}=\int_{t_0}^{t_1}|v(t)|\dd t
\end{equation}
とする。速さは非負，速度には向きに対応する符号がある。この違いは，面積と符号付きの定積分の違いと同じである。

速度が零になる時刻は，進む向きが変わる候補である。ただし零になるだけでは反転したとは限らない。たとえば $v(t)=(t-1)^2$ は $t=1$ で零になるが，前後で非負である。道のりを計算するときは，零点で区間を分けたうえで，各区間の符号を確かめる。

\begin{example}{変位と道のりを分けて求める}
時刻 $0\leq t\leq4$ において，速度が $v(t)=(t-1)(t-3)$ で与えられる。初期位置は $s(0)=0$ とする。$t=4$ の位置，道のり，平均速度，平均の速さを求めよ。
\end{example}

速度を積分して初期条件を使うと，
\[
s(t)=\frac{t^3}{3}-2t^2+3t.
\]
したがって $s(4)=\frac43$ である。道のりでは速度の符号を調べる。$v$ は $0\leq t<1$ で正，$1<t<3$ で負，$3<t\leq4$ で正である。折り返す位置は
\[
s(1)=\frac43,\qquad s(3)=0
\]
なので，道のりは
\[
\left|\frac43-0\right|+\left|0-\frac43\right|
 +\left|\frac43-0\right|=4.
\]
これは $\int_0^1v\dd t-\int_1^3v\dd t+\int_3^4v\dd t$ と計算しても同じである。全時間は $4$ だから，平均速度は $\cfrac{\cfrac43}{4}=\frac13$，平均の速さは $\frac44=1$ となる。

加速度は $a(t)=2t-4$ である。加速度が正でも，いつも速さが増えるわけではない。$v\ne0$ の時刻では
\[
\frac{\dd}{\dd t}|v(t)|=
\begin{cases}
a(t),&v(t)>0,\\
-a(t),&v(t)<0
\end{cases}
\]
だから，速度と加速度が同符号なら速さが増え，異符号なら減る。この例では $1<t<2$ と $3<t<4$ で速さが増える。負の向きに進みながら速くなることもあるのである。

平面上の運動 $\bigl(x(t),y(t)\bigr)$ では，速度は二成分 $\bigl(x'(t),y'(t)\bigr)$ を持ち，速さは $\sqrt{x'(t)^2+y'(t)^2}$ となる。その積分が道のりになることは，式\eqref{式:媒介曲線の弧長}そのものである。運動では同じ場所を再び通った分も道のりに含むので，曲線を重複してたどることは誤りではなく，数えるべき内容になる。

\section{連続的に変わる量の平均}\label{節:関数の平均}

\sectionref{節:定積分性質}で導入した関数の平均値を，ここでは温度のような連続的に変わる量へ使おう。まず，平均を積分で表す理由を振り返る。数 $u_1,\ldots,u_n$ の平均は和を個数で割ったものである。関数 $f$ の区間 $[a,b]$ における平均も，等間隔に値を取り，その平均の極限として考えられる。$\Delta x=\frac{b-a}{n}$ とすれば，
\[
\frac1n\sum_{k=1}^n f(a+k\Delta x)
 =\frac{1}{b-a}\sum_{k=1}^n f(a+k\Delta x)\Delta x.
\]
右辺の和はリーマン和なので，$f$ が連続なら極限は
\begin{equation}\label{式:関数の平均値}
\overline f=\frac{1}{b-a}\int_a^b f(x)\dd x
\end{equation}
となる。$f\geq0$ の場合，底辺 $b-a$，高さ $\overline f$ の長方形は，曲線の下と同じ面積になる。$f$ に負の値があっても，平均の定義は同じ式である。

\begin{example}{変化する温度の平均}
時刻 $0\leq t\leq12$ の温度を $T(t)=20+6\sin\frac{\pi t}{12}$ とする。この時間帯の平均温度を求めよ。
\end{example}

始めと終わりの温度はともに $20$ だが，途中の温度は $20$ より高い。両端の平均だけでは時間全体の平均にならない。定義に従って
\[
\begin{aligned}
\overline T
 &=\frac1{12}\int_0^{12}\left(20+6\sin\frac{\pi t}{12}\right)\dd t\\
 &=20+\frac12\left[-\frac{12}{\pi}\cos\frac{\pi t}{12}\right]_0^{12}
 =20+\frac{12}{\pi}.
\end{aligned}
\]
全時間を均等に数えて平均を取っていることが，$\frac1{12}$ に表れている。

連続な $f$ の最小値を $m$，最大値を $M$ とすると，$m\leq f(x)\leq M$ を積分して
\[
m\leq\overline f\leq M
\]
を得る。中間値の定理から，ある $c\in[a,b]$ で $f(c)=\overline f$ となる。これを積分の平均値の定理という。平均が実際のどこかの値と一致することは分かるが，一般にその点が区間の中央とは限らず，一つとは限らない。

均等に数えるのでない平均には，重みが必要になる。非負の連続関数 $w$ に対し $\int_a^b w(x)\dd x>0$ なら，重み付き平均は
\[
\frac{\int_a^b f(x)w(x)\dd x}{\int_a^b w(x)\dd x}
\]
である。分母は重みの総量であり，$w=1$ の場合が式\eqref{式:関数の平均値}になる。次に扱う重心は，位置の座標を面積や長さで重み付けして平均したものである。

\section{発展：面積の重心と線の重心}\label{節:面積重心と線重心}

密度が一様な薄い板を考える。板を小さな部分に分け，それぞれの面積を $\Delta A_i$，位置を $(x_i,y_i)$ と近似すると，重心の座標は面積を重みにした平均
\[
\overline x\approx\frac{\sum_i x_i\Delta A_i}{\sum_i\Delta A_i},\qquad
\overline y\approx\frac{\sum_i y_i\Delta A_i}{\sum_i\Delta A_i}
\]
で表される。密度が一様なら各部分の質量が面積に比例し，共通の密度が分子と分母で消えるからである。

領域が $a\leq x\leq b$，$g(x)\leq y\leq f(x)$ で表され，面積 $A>0$ とする。縦の細い長方形の面積はほぼ $(f-g)\Delta x$ で，その中心の座標は $\bigl(x,\frac{f+g}{2}\bigr)$ である。分割を細かくすると
\begin{equation}\label{式:面積の重心}
\begin{aligned}
A&=\int_a^b(f-g)\dd x,\\
\overline x&=\frac1A\int_a^b x(f-g)\dd x,\\
\overline y&=\frac1{2A}\int_a^b(f^2-g^2)\dd x
\end{aligned}
\end{equation}
を得る。最後の式は $\frac{f+g}{2}(f-g)=\frac12(f^2-g^2)$ から出ている。$\overline y$ を計算するときに上端の高さ $f$ だけを掛けてはいけない。帯全体の重心は上下の中央にある。

曲線に沿って曲げた一様な針金では，質量が比例するのは面積ではなく長さである。長さを $L>0$ とすると，重心は
\begin{equation}\label{式:線の重心}
\overline x=\frac1L\int_C x\dd s,\qquad
\overline y=\frac1L\int_C y\dd s
\end{equation}
となる。ここで $\dd s$ は曲線の小さな長さを表す記号で，媒介表示なら
\[
\int_C y\dd s
 =\int_\alpha^\beta y(t)\sqrt{x'(t)^2+y'(t)^2}\dd t
\]
という普通の定積分で計算する。一様な針金では，一度だけたどる区間を選ぶ。

半径 $a$ の上半円板と，その円弧だけの針金を比較する。どちらも左右対称なので $\overline x=0$ である。半円板では $A=\frac{\pi a^2}{2}$，$g=0$，$f(x)=\sqrt{a^2-x^2}$ だから，
\[
\int_A y\dd A
 =\frac12\int_{-a}^a(a^2-x^2)\dd x
 =\frac{2a^3}{3}.
\]
左辺は面積で重み付けした $y$ 座標の和を表す記号であり，この場合の計算方法が右辺である。したがって，半円板の重心の高さは
\begin{equation}\label{式:半円板の重心}
\overline y_{\mathrm{面}}=\frac{4a}{3\pi}.
\end{equation}
円弧のほうは $x=a\cos t$，$y=a\sin t$，$0\leq t\leq\pi$ と表すと，$\dd s=a\dd t$，$L=\pi a$ だから，
\begin{equation}\label{式:半円弧の重心}
\overline y_{\mathrm{線}}
 =\frac{1}{\pi a}\int_0^\pi(a\sin t)a\dd t
 =\frac{2a}{\pi}.
\end{equation}
円弧だけのほうが重心は高い。円板には直径の近くにも面積があるが，円弧だけの針金にはその部分がないためである。

\section{発展：パップス＝ギュルダンの二つの定理}\label{節:パップスの定理}

回転体の体積を円筒殻の式で求めると，距離に面積を掛けた積分が現れる。一方，重心も座標に面積を掛けた積分で求めた。この共通する積分を使うと，回転体の体積を面積と重心で表せる。

まず回転軸を $x$ 軸とし，領域がその上側にある場合を考える。高さ $y$ の横断面の長さを $h(y)$ とする。横の細い帯を回すと半径 $y$，軸に沿う長さ $h(y)$ の円筒殻になるので，
\[
V=2\pi\int y h(y)\dd y.
\]
他方，面積 $A$ と重心の高さ $\overline y$ は
\[
A=\int h(y)\dd y,\qquad
A\overline y=\int y h(y)\dd y
\]
を満たす。よって $V=2\pi\overline y A$ である。積分区間は領域に含まれる高さ全体であり，断面が複数の区間に分かれるなら，$h(y)$ はそれらの長さの和とする。

\begin{theorem}{パップス＝ギュルダンの体積に関する定理}
一様な平面図形の面積を $A$，その面積重心と，図形と同じ平面内にある回転軸との距離を $R$ とする。図形が軸の片側にあり，回転によって内部を重複して数えない場合，できる回転体の体積は
\[
V=2\pi R A
\]
である。すなわち，面積に，重心が一回転するときに進む円周の長さを掛ければよい。
\end{theorem}

ここでは境界が区分的に滑らかで，通常の断面積による積分ができる図形を考えている。軸が境界に触れても内部を横切らなければ，上の帯による証明を使える。軸が内部を横切る場合には，軸の両側が同じ場所を回転して重なり得る。重心が軸上にあるから体積が零になる，という使い方はできない。

\subsection{回転面の面積と線重心}

回転軸と同じ平面内にある曲線を考える。この曲線を回したときにできる面の面積では，面積重心ではなく線重心が現れる。その理由を，回転面の小さな部分から考える。軸の片側にあり，軸からの距離が $r_1,r_2$ の二点を結ぶ短い線分を回すと，一般には円すい台の側面になる。軸に平行な線分は円筒面，垂直な線分は円環または円板を作る。いずれの場合も，線分の長さを $\Delta s$ とすると面積は
\[
\pi(r_1+r_2)\Delta s
\]
である。円すい台の場合は，大きい円すいと切り取った円すいの側面積の差を，相似比で整理するとこの式を得る。円筒面と円環の場合も，それぞれの面積公式から確かめられる。線分を十分短くすれば，$r_1,r_2$ はほぼ同じ $r$ なので，これは $2\pi r\Delta s$ になる。

回転軸と同じ平面内の曲線が区分的に連続微分可能で，回転によって同じ面を重複して作らないとき，折れ線近似の極限として回転面の面積は
\begin{equation}\label{式:回転面の面積}
S=2\pi\int_C r\dd s
\end{equation}
になる。ここでは，この極限の正当化の詳細は省く。たとえば $y=f(x)\geq0$ を $x$ 軸のまわりに回すなら
\[
S=2\pi\int_a^b f(x)\sqrt{1+f'(x)^2}\dd x.
\]
ここで得られるのは，指定した曲線が作る面の面積である。立体の両端を塞ぐ円板などは，その曲線に含めていなければ別に数える。

\begin{theorem}{パップス＝ギュルダンの面積に関する定理}
一様な平面曲線の長さを $L$，その線重心と，曲線と同じ平面内にある回転軸との距離を $R$ とする。曲線が軸の片側にあり，一回転によって同じ面を重複して作らない場合，曲線が作る回転面の面積は
\[
S=2\pi R L
\]
である。
\end{theorem}

実際，軸を $x$ 軸，曲線を $y\geq0$ に取れば $r=y$ であり，線重心の定義から $\int_C y\dd s=L\overline y=LR$ となるので，式\eqref{式:回転面の面積}から従う。回転軸の位置や方向が違っていても，軸に合わせて座標を取り直せば同じ議論になる。

二つの定理で掛ける円周は，図形の端や幾何学的な中心の円周ではなく，それぞれの重心が進む円周である。また，体積の定理には面積 $A$ と面積重心，面の定理には長さ $L$ と線重心を対応させる。

\begin{example}{半円板を外側の軸のまわりに回す}
$a>0$，$b>0$ とする。上半円板 $x^2+y^2\leq a^2$，$y\geq0$ を，直線 $y=-b$ のまわりに一回転してできる立体の体積と全表面積を求めよ。
\end{example}

図\ref{図:半円の二つの重心}のように，回転軸は図形の外にあり，重心の高さに $b$ を加えたものが回転半径になる。半円板の面積重心を使うと，
\[
\begin{aligned}
V&=2\pi\left(b+\frac{4a}{3\pi}\right)\frac{\pi a^2}{2}\\
 &=\pi^2a^2b+\frac{4\pi a^3}{3}.
\end{aligned}
\]
断面で確かめるなら，外半径 $b+\sqrt{a^2-x^2}$，内半径 $b$ の円環である。したがって
\[
\begin{aligned}
V&=\pi\int_{-a}^a
\left\{\left(b+\sqrt{a^2-x^2}\right)^2-b^2\right\}\dd x\\
 &=\pi\int_{-a}^a(a^2-x^2)\dd x
  +2\pi b\int_{-a}^a\sqrt{a^2-x^2}\dd x,
\end{aligned}
\]
となり，第一項が $\frac{4\pi a^3}{3}$，第二項が $\pi^2a^2b$ で一致する。

\begin{figure}[htbp]
\centering
\begin{tikzpicture}[scale=1.25,>=Stealth]
\fill[blue!10](-2,0)arc[start angle=180,end angle=0,radius=2]--cycle;
\draw[blue!75!black,thick](2,0)arc[start angle=0,end angle=180,radius=2];
\draw[thick](-2,0)--(2,0);
\draw[->](-2.5,0)--(2.6,0)node[right]{$x$};
\draw[->](0,-1.5)--(0,2.45)node[above]{$y$};
\draw[orange!80!black,thick](-2.5,-1)--(2.5,-1)node[right]{$y=-b$};
\fill[blue!75!black](0,0.8488)circle(0.045);
\fill[red!75!black](0,1.2732)circle(0.045);
\draw[blue!75!black](0,0.8488)--(2.5,0.6)node[right]{面積重心 $\frac{4a}{3\pi}$};
\draw[red!75!black](0,1.2732)--(2.5,1.6)node[right]{線重心 $\frac{2a}{\pi}$};
\draw[<->](-2.25,-1)--(-2.25,0)node[midway,left]{$b$};
\node[below]at(-2,0){$-a$};\node[below]at(2,0){$a$};
\end{tikzpicture}
\caption{同じ半円でも，板と円弧では重心が違う。回転半径は軸から測る。}\label{図:半円の二つの重心}
\end{figure}

表面積では，境界を円弧と直径に分ける。円弧が作る面は，円弧の長さ $\pi a$ と線重心の高さ $\frac{2a}{\pi}$ を使って
\[
S_{\mathrm{円弧}}
 =2\pi\left(b+\frac{2a}{\pi}\right)\pi a
 =2\pi^2ab+4\pi a^2.
\]
これだけでは全表面積にならない。直径 $-a\leq x\leq a$，$y=0$ も，回転すると半径 $b$，軸方向の長さ $2a$ の円筒面を作る。その面積は
\[
S_{\mathrm{直径}}=2\pi b\cdot2a=4\pi ab.
\]
両者は端の円周で接するが，面積のある部分では重ならない。よって全表面積は
\[
S=2\pi^2ab+4\pi a^2+4\pi ab.
\]
$x=\pm a$ では外半径と内半径がともに $b$ となり，厚みが零なので，さらに端の円板を加える必要はない。どの境界がどの面を作るかを確かめた結果として，全表面積が決まる。

面積，体積，長さ，道のり，平均はいずれも積分で表せるが，同じ被積分関数を使うわけではない。細い帯なら高さと幅，薄い立体なら断面積と厚さ，短い曲線なら二方向の変化から作った距離が必要になる。公式を選ぶ前に，細かく分けた一つ分が何を表すかを言えるようにすると，符号や重複の誤りに気づきやすくなる。

\section*{章末問題}
\begin{enumerate}[label=\textbf{\thechapter.\arabic*},leftmargin=*]
\item\label{問:面積と円環}
$y=x^2$ と $y=3x$ に囲まれた部分の面積を求めよ。また，この部分を $x$ 軸のまわりに回転させた立体の体積を求めよ。
\item\label{問:横帯の面積}
$x=y^2$ と $x=3-2y$ に囲まれた部分の面積を，$y$ を積分変数として求めよ。
\item\label{問:外半径の交代}
$0\leq x\leq2$，$-1\leq y\leq x$ の領域を $x$ 軸のまわりに回した立体の体積を求めよ。外半径が変わる位置を明示せよ。
\item\label{問:円筒と円板}
$0\leq x\leq1$，$x^2\leq y\leq1$ の領域を $y$ 軸のまわりに回した立体の体積を，円筒殻と円板の二通りで求めよ。
\item\label{問:弧長の平方完成}
$y=\frac13(x^2+2)^{\frac32}$，$0\leq x\leq1$ の曲線の長さを求めよ。
\item\label{問:道のりと平均速度}
$0\leq t\leq3$ における速度が $v(t)=3(t-1)(t-2)$，初期位置が $s(0)=0$ の運動について，終点の位置，道のり，平均速度，平均の速さを求めよ。
\item\label{問:平均値を取る点}
$f(x)=1+x^2$ の $[0,2]$ における平均値を求め，$f(c)$ がその平均値と等しくなる $c\in[0,2]$ を求めよ。
\item\label{問:三角形と両パップス}
頂点 $(0,0),(2,0),(2,1)$ の三角形を考える。面積重心の $y$ 座標を積分で求め，パップス＝ギュルダンの定理により，三角形を $x$ 軸のまわりに回した立体の体積を求めよ。また，境界の各線分の線重心を使って，できる立体の全表面積を求めよ。回転軸にある線分の寄与も説明せよ。
\end{enumerate}

\smallskip\noindent 解答は\pageref{解答:第8章}ページ。
